%% =====================================================================
%%  structures/informe/structure.tex  --  Esqueleto del INFORME
%% ---------------------------------------------------------------------
%%  Informe tecnico de despacho:
%%
%%     INFORME N° 001-2026/ABC/CR
%%     A / DE / ASUNTO / FECHA
%%     I. ANTECEDENTES   II. ANALISIS   III. CONCLUSIONES ...
%%
%%  Las secciones se numeran solas en romanos: \seccion{ANTECEDENTES}.
%% =====================================================================
\ifcsname despachoEstructuraInforme\endcsname
  \expandafter\endinput
\fi
\expandafter\def\csname despachoEstructuraInforme\endcsname{cargada}

\newcommand{\numerodocumento}{%
  \begin{center}
    \ifmetadef{numerocompleto}%
      {\underline{\seriesemi\MakeUppercase{\meta{numerocompleto}}}}%
      {\underline{\seriesemi\MakeUppercase{\meta{tipodoc}}\ N°~\meta{numero}-\meta{anio}/\meta{siglas}}}%
  \end{center}
  \vspace{0.3\baselineskip}%
}

%%  \cabecerainforme[referencia]{a}{de}{asunto}
%%  La REFERENCIA es OPCIONAL (informes que responden a un pedido). El
%%  \\ de la fila viaja dentro de \despachorefinf para no quedar
%%  adyacente a un \fi dentro del tabular (gotcha de \@arraycr).
\newcommand{\cabecerainforme}[4][]{%
  \def\despachorefinf{}%
  \if\relax\detokenize{#1}\relax\else
    \def\despachorefinf{\textbf{REF.} & #1 \\}%
  \fi
  \noindent
  \begin{tabular}{@{}p{1.9cm}@{~:~}p{\dimexpr\textwidth-2.6cm\relax}@{}}
    \textbf{A}      & #2 \\
    \textbf{DE}     & #3 \\
    \textbf{ASUNTO} & #4 \\
    \despachorefinf
    \textbf{FECHA}  & \meta{lugar}, \ifmetadef{fecha}{\meta{fecha}}{\today} \\
  \end{tabular}
  \par\vspace{2pt}
  \noindent{\color{neutral500}\rule{\textwidth}{0.5pt}}%
  \par\espacioantescuerpo{0.8}% aire hasta el primer parrafo (ajustable, 2026-10-01)
  \setmeta{asunto}{#4}%
}

%% --- Secciones en romanos (I., II., ...) ------------------------------
\newcounter{seccioninforme}
\newcommand{\seccion}[1]{%
  \stepcounter{seccioninforme}%
  \vspace{0.9\baselineskip}%
  \noindent{\fuentedisplaysemi\color{colortitulos}%
    \Roman{seccioninforme}.\ \MakeUppercase{#1}}%
  \par\vspace{0.35\baselineskip}\noindent\ignorespaces
}

\newcommand{\atentamente}[1][Es cuanto informo a usted para los fines pertinentes.]{%
  \vspace{0.4\baselineskip}%
  \noindent #1\par
}

%% (\bloquefirma, \adjunto y \concopia viven en core/lib/partes-comunes.tex)
